% =================================
% Experiments + Results 
% =================================
\section{Experiments}
\label{sec:exp_results}

%% NOTE: Ebasa's original scripts
% Our experiments are designed to test the capabilities of the learned policy by answering the following four questions:
% (1) \textbf{Performance:} How accurately can the robot push the object to achieve the goal pose?
% (2) \textbf{Generalization:} does performance hold across large distributions of object geometry and physical properties?
% (3) \textbf{Robustness:} how sensitive is performance to initializations?
% (4) \textbf{Behavior:} how does the model perform for two different contact modes of manipulation i.e. \textit{whole-body} vs. \textit{arm-mediated}?
% \Xun{Remove this paragraph}
% In this section we will define the experimental setup, task, training and testing distributions, along with the evaluation protocol and metrics. 

% Subsections following this could be: A) Non-prehensile Manipulation Task and Contact Modes (unless this comes in the problem statement); B) Task Success Metrics; C) Evaluation Protocol for Testing Generalization; D) Evaluation Protocol for Testing Robustness.

To assess the effectiveness of the trained policy in pushing the object to a desired goal pose, and to evaluate its capacity to generalize to out-of-distribution conditions (e.g., previously unseen object geometries or mass properties), we conduct a series of controlled experiments in simulation. Section \ref{subsec:exp_setup} first outlines the experimental environment implemented in IsaacLab \cite{Baginski2025IsaacLab}, including the task configuration for pushing the target object to a specified goal pose. Section \ref{subsec:task_protocol} subsequently details the experimental variations adopted to assess the robustness and generalization of the trained policy on out-of-distribution conditions. Finally, Section \ref{subsec:eval_metrics} describes the evaluation metrics used to quantify policy performance and to compare the proposed approach against the baseline method.

% \kar{Our experiments are designed to test the capabilities of the learned policy by answering the following four questions:
% (1) \textbf{Performance:} How accurately can the robot push the object to achieve the goal pose?
% (2) \textbf{Generalization:} does performance hold across large distributions of object geometry and physical properties?
% (3) \textbf{Robustness:} how sensitive is performance to initializations?
% (4) \textbf{Behavior:} how does the model perform for two different contact modes of manipulation i.e. \textit{whole-body} vs. \textit{arm-mediated}?
% }

% \kar{In this section we will define the experimental setup, task, training and testing distributions, along with the evaluation protocol and metrics. 
% }

% \kar{Subsections following this could be: A) Non-prehensile Manipulation Task and Contact Modes (unless this comes in the problem statement); B) Task Success Metrics; C) Evaluation Protocol for Testing Generalization; D) Evaluation Protocol for Testing Robustness.}

% ================================
\subsection{Experimental Setup}
\label{subsec:exp_setup}

Consistent with the training procedure, the benchmarking environment is implemented in IsaacLab \cite{Baginski2025IsaacLab}. At the beginning of each episode, the robot, the target object, and the goal pose are initialized with random poses within an $8\mathrm{m} \times 8\mathrm{m}$ planar arena. The ground plane is assigned a friction coefficient $\mu$, sampled independently to introduce variability in contact dynamics (see (TODO: a figure on the arena)). The set of object categories evaluated in our experiments is summarized in Table~\ref{tab:object_sets}.

The task objective is to control the robot to push the initialized object to the specified goal pose, either using its base or manipulator, while maintaining dynamic stability and satisfying predefined contact-safety constraints. An episode is considered unsuccessful and terminated if any constraint is violated, including object toppling or joint torque saturation. 

% =========================
% TABLE: setup summary 
% =========================
\begin{table}[t]
\centering
\caption{\textbf{Evaluation setup summary.}}
\label{tab:setup_summary}
\small
\setlength{\tabcolsep}{4pt}
\renewcommand{\arraystretch}{1.05}
\begin{tabular}{@{}p{2.9cm}cc@{}}
\toprule
 & \textbf{WB} & \textbf{AM} \\
\midrule
Robot & [Spot, ANYmal-C]  & [Spot + arm] \\
High-level action dim. & [3] & [12] \\
Policy obs dim. & [25] & [72] \\
Episode horizon (s) & [30] & [30] \\
Object families & [box,cyl,chair] & [box,cyl,chair] \\
Permitted contacts & base/body links & arm/EE links \\
Parallel envs & \multicolumn{2}{c}{[4096]} \\
% Rollouts per suite ($N$) & \multicolumn{2}{c}{[1000]} \\
% Seeds ($S$) & \multicolumn{2}{c}{[3]} \\
\bottomrule
\end{tabular}
\vspace{-0.75em}
\end{table}

% ============================================================
\subsection{In-Distribution (ID) \& Out-of-Distribution (OOD) Test Suites}
\label{subsec:task_protocol}


% \paragraph{Task}
% NOTE: Ebasa's original scripts
% Each episode samples an object instance and a planar goal pose. The objective is to move the object from its initial pose to the goal pose while maintaining stability and respecting contact-safety constraints. Episodes terminate on timeout, object failure (e.g., topple), or robot/safety violations.


% At the time of spawning, the attributes of the objects are all randomly initialized. By comparing with the training dataset, the spawned objects can be classified into two types: \textit{in-distribution (ID)} or \textit{Out-of-Distribution (OOD)}. Here, \textit{in-distribution (ID)} refers to the objects with a random initialization distribution that matches the training distribution, while \textit{out-of-distribution (OOD)} are the ones with a random initialization distribution that deliberately shifts the training distribution to diagnose generalization.
% \textit{OOD} objects are initialized under conditions that lie outside the training distribution, e.g., novel object geometry or physics parameters beyond the training randomization ranges.
% We isolate geometry shifts, physics shifts, and difficult initializations: 
% \begin{enumerate}
%     \item \textit{ID:} object families and parameter ranges seen during training.
%     \item \textit{OOD-Geometry (OOD-G):} novel shapes and aspect ratios (e.g., long/flat cuboids, tall/slender objects, non-convex variants).
%     \item \textit{OOD-Physics (OOD-P):} friction/mass/CoM/inertia variations outside the training range.
%     \item \textit{Stress:} harder initializations (tight clearances, random yaw, near-boundary goals, cluttered starts).
% \end{enumerate}

% For the object listed in (Table~\ref{tab:object_sets}) and used for our experiments, we detail our ID/OOD test suites in Table~\ref{tab:domain_rand }, according to whether geometry/physics parameters fall inside or outside the train-time randomization ranges 

At episode initialization, all object attributes are randomly sampled. Relative to the training distribution, the instantiated objects are categorized as either in-distribution (ID) or out-of-distribution (OOD). The ID setting corresponds to objects whose geometric and physical parameters are drawn from the same randomization ranges employed during training. In contrast, OOD instances are generated by systematically shifting the sampling distribution to evaluate the policy’s generalization under distributional mismatch.

Specifically, OOD objects are initialized with properties that lie outside the support of the training-time domain randomization, such as previously unseen geometries or physical parameters that exceed the nominal training bounds. To disentangle the sources of distributional shift, we separately examine geometric variations and physical parameter variations:
\begin{enumerate}
    \item \textit{ID: } Object families and parameter ranges observed during training.
    \item \textit{OOD-Geometry (OOD-G): }Novel shapes and aspect ratios (e.g., elongated or flattened cuboids, tall and slender objects, and non-convex variants).
    \item \textit{OOD-Physics (OOD-P):} Variations in friction, mass, center of mass, or inertia that fall outside the training randomization ranges.
\end{enumerate}


% =========================
% TABLE II: Object taxonomy + evaluation sets
% =========================
\begin{table}[t]
\caption{\textbf{Object taxonomy and evaluation sets.} }
\label{tab:object_sets}
\renewcommand{\arraystretch}{1.05}
\setlength{\tabcolsep}{4pt}
\small
\begin{tabular}{@{}p{0.6cm}p{2.4cm}p{1.35cm}p{0.5cm}p{2.8cm}@{}}
\toprule
\textbf{Cat.} & \textbf{Family} & \textbf{Type} & \textbf{\#} & \textbf{Suites} \\
\midrule
\multirow{2}{*}{C1} & Cuboid & parametric & 500 & ID, OOD-G, OOD-P \\
                    & Cylinder & parametric & 500 & ID, OOD-G, OOD-P \\
\midrule
\multirow{2}{*}{C2} & Chair & rigid mesh & 9 & ID, OOD-G, Stress \\
                    & Stool / table & rigid mesh & - & OOD-G, Stress \\
\midrule
\multirow{2}{*}{C3} & Chair-with-casters & articulated & - & OOD-Art, Stress \\
                    & Cart / dolly & articulated & - & OOD-Art, Stress \\
\bottomrule
\end{tabular}
\vspace{-0.6em}
\end{table}



% % % ============================================================
% % \paragraph{Domain Randomization}
% % When initializing our evaluation scenario, we randomize the domain according to the statistics in \ref{tab:domain_rand}. Notably, "Train" means the statistics are adopted in the training process, while "OOD" (Out-of-Distribution) refers to the statistics unseen in the original training process.To make ID/OOD definitions explicit, we report train-time randomization ranges and the expanded OOD ranges as seen in Table \ref{tab:domain_rand}. 

% =========================
% TABLE I: Domain randomization summary
% =========================
\begin{table}[t]
\centering
\caption{\textit{Domain randomization}}
\label{tab:domain_rand}
\renewcommand{\arraystretch}{1.05}
\setlength{\tabcolsep}{4pt}
\small
\begin{tabular}{@{}p{4.05cm}p{2.05cm}p{2.05cm}@{}}
\toprule
\textbf{Factor} & \textbf{Train}  & \textbf{OOD} \\
\midrule
Object mass $m$ (kg) & $[2.0,8.0]$  & $[2,18]$ \\
Static friction $\mu$ & $[0.6,1.2]$  & $[0.1,0.6]$ \\
Dynamic friction $\mu_d$ & $[0.48, 0.96]$ & $[0.08, 0.48]$ \\
CoM offset $\|\Delta c\|$ (m) & $\le 0.05$ & $\le 0.20$ \\
% Inertia tilt (deg) & $[-15,15]$ & $[-25,25]$ \\
% Drag / damping & $[0.05,1.0]$  & $[0.0,1.5]$ \\
\midrule
Object Height $h$ & $[0.4,1.0]$  & $[0.4,1.8]$ \\
Object width $w$ (m) & $[0.25, 0.75]$ & $[0.25, 1.25]$ \\
Object length $d$ (m) & $[0.25, 0.75]$ & $[0.25, 1.25]$ \\
% Aspect ratios & bounded  & expanded \\
% Init yaw & $[-\pi,\pi]$  & same \\
\midrule
Spawn region $(x,y)$ (m) & $x,y \in [-4,4]$ & same \\
Init yaw $\psi$ (rad) & $[-\pi,\pi]$ & same \\
\midrule
Goal distance $\|p_g-p_o\|$ (m) & $[1.5,2.5]$  & $[0.8,3.5]$ \\
Robot--obj dist $\|p_o-p_b\|$ (m) & $[2.0,3.2]$  & same \\
Episode time limit $T_{\max}$ (s) & 30 & 30 \\
% Tight clearances & no  & yes \\
\midrule
% External pushes & on/off  & stronger \\
% Sensor noise & on/off  & stronger \\
Terrain & plane & plane/rough \\
\bottomrule
\end{tabular}
% \vspace{-0.6em}
\end{table}

For the object categories listed in Table~\ref{tab:object_sets} and evaluated in our experiments, the corresponding ID and OOD test suites are summarized in Table~\ref{tab:domain_rand}, organized according to whether the geometric and physical parameters fall within or beyond the training-time randomization ranges.

%%%%%%%%%%%%%%%%%%%%%%%%%%%%%%%%%%%%%%%%%%%%%%
%%%%  Training Results from Spot (no arm) %%%%
%%%%%%%%%%%%%%%%%%%%%%%%%%%%%%%%%%%%%%%%%%%%%%
% \begin{table}[t]
% \centering
% \caption{\textbf{Domain randomization} Train-time randomization factors.}
% \label{tab:domain_rand}
% \renewcommand{\arraystretch}{1.05}
% \setlength{\tabcolsep}{4pt}
% \small
% \begin{tabular}{@{}p{3.6cm}p{2.2cm}@{}}
% \toprule
% \textbf{Factor} & \textbf{Train} \\
% \midrule
% Object mass $m$ & [6, 18] \\
% Friction $\mu$ & [0.2, 0.4] \\
% CoM offset $\|\Delta c\|$ & [10, 40] \\
% Inertia tilt & [-15, 15] \\
% Drag / damping & [0.05, 1] \\
% Geom. scale $s$ & -- \\
% Aspect ratios & -- \\
% Init yaw & [0, 6.28319] \\
% Goal distance & goal\_radius=0.35, resample=[1e+09, 1e+09] \\
% Robot--obj distance & reset\_base pose x=[-4, 4], y=[-4, 4] \\
% Tight clearances & -- \\
% External pushes & -- \\
% Sensor noise & -- \\
% Terrain & type=plane; ground\_friction\_range=[0.3, 1.2]; ground\_friction\_pattern=random \\
% \bottomrule
% \end{tabular}
% \vspace{-0.6em}
% \end{table}

\subsection{Baseline \& Evaluation Metrics}
\label{subsec:eval_metrics}
% \paragraph{Baseline Methods} We compare two \emph{contact modalities} under an identical goal-conditioned pushing task:
% \begin{itemize}
%     \item \textit{Whole-body pushing} (WB): the robot primarily uses base/body contacts to push and steer the object.
%     \item \textit{Arm-mediated pushing} (AM): the robot uses end-effector/arm contacts to push while coordinating locomotion.
% \end{itemize}
% Both modalities share the same goal specification and evaluation metrics, but differ in action/observation interfaces and permissible contact sets (Sec.~\ref{sec:methods}).
% \paragraph{Metrics}
% We report: 
% (i) success at two tolerances, where the Object reachs the goal with (15cm, 15deg) and stricter metrics where the object reachs the goal under (10cm,10deg) difference;
% (ii) terminal position and yaw errors; 
% (iii) failure attribution by threshold component (pos-only / yaw-only / both);
% (iv) stability (upright/tilt, angular velocity); (v) safety/contact (undesired contacts, illegal contacts, non-arm contact share);
% and (vi) efficiency (time-to-success, robot path length, push efficiency).


% \paragraph{Tolerance-robustness and AUC.}
% To understand performance beyond a single threshold, we evaluate success on a grid of tolerances $(e_{\text{pos}}, e_{\psi})$.
% We summarize \emph{robustness} by the \textit{area under the curve (AUC)}: a single number that measures how quickly success improves as tolerances are relaxed.
% Concretely, we compute a normalized AUC over the plotted range:
% \[
% \mathrm{AUC} \;=\; \frac{1}{|\mathcal{G}|}\sum_{(p,\psi)\in\mathcal{G}} \mathrm{Success}(p,\psi),
% \]
% where $\mathcal{G}$ is the tolerance grid. Higher AUC indicates a policy that succeeds across a wider range of strict-to-lenient criteria (i.e., more robust alignment).


\paragraph{Baseline Methods}
We compare two distinct \emph{contact modalities} under an identical goal-conditioned pushing task:
\begin{itemize}
\item \textit{Whole-body pushing} (WB), in which the robot primarily exploits base or body contacts to propel and steer the object;
\item \textit{Arm-mediated pushing} (AM), in which the robot employs end-effector or arm contacts while coordinating locomotion to achieve object displacement.
\end{itemize}
Both modalities share the same goal specification and evaluation protocol, ensuring a controlled comparison. They differ, however, in their action and observation interfaces, as well as in the set of admissible contacts (see Sec.~\ref{sec:methods}).

\paragraph{Metrics}
When the execution of the policy reaches the time limit, we report the following performance measures based on the final pose of the object (x, y, $\theta$) and the goal pose:
\begin{enumerate}
    \item task success under two tolerance thresholds: \textbf{Success@(15, 15)}(If the positional error is smaller than 15 cm and the angular error is smaller than 15$^\circ$) and a stricter \textbf{Success@(10, 10)}(If the positional error is smaller than 10 cm and the angular error is smaller than 10$^\circ$); 
    %\item terminal position and yaw errors;
    \item failure attribution categorized by threshold component (position-only, yaw-only, or both);
    % \item stability indicators, including uprightness/tilt and angular velocity;
    % \item safety and contact metrics, such as undesired or illegal contacts and the proportion of non-arm contacts;
    \item efficiency metrics, including time-to-success, robot path length, and push efficiency.
\end{enumerate}


\paragraph{Tolerance Robustness and AUC}
To characterize performance beyond a single tolerance threshold, we evaluate success over a grid of position and orientation tolerances $(e_{\text{pos}}, e_{\psi})$. Robustness is summarized via the \emph{area under the curve} (AUC), which provides a scalar measure of how rapidly success rates increase as tolerances are relaxed. Concretely, we compute a normalized AUC over the evaluated tolerance grid:

\[
\mathrm{AUC} \;=\; \frac{1}{|\mathcal{G}|}\sum_{(p,\psi)\in\mathcal{G}} \mathrm{Success}(p,\psi),
\]

where $\mathcal{G}$ denotes the set of tolerance pairs. A higher AUC indicates that the policy achieves reliable performance across a broader spectrum of strict-to-lenient evaluation criteria, reflecting greater robustness in goal alignment.